\documentclass[class=llncs, crop=false]{standalone}

\usepackage{unicode-math}
\usepackage{newcomputermodern}
\usepackage{microtype}
\usepackage{newunicodechar}
\newunicodechar{⤳}{\ensuremath{⤳}}
\newunicodechar{↪}{\ensuremath{↪}}
\newunicodechar{≔}{\ensuremath{≔}}
\newunicodechar{⊢}{\ensuremath{⊢}}

% \usepackage{amsthm}
\usepackage{xcolor-material}
\usepackage{etoolbox}

%% ── Draft toggle ─────────────────────────────────────────────────
%% Comment out the next line for submission.
% \newcommand{\darkmode}{}

\ifdefined\darkmode
	\usepackage{showframe}
	\renewcommand\ShowFrameLinethickness{0.2pt}
	\renewcommand*\ShowFrameColor{\color{MaterialGrey600}}
	\usepackage{pagecolor}
	\pagecolor{MaterialGrey900}
	\color{MaterialGrey100}
\fi

\usepackage{tikz}
\usetikzlibrary{fit}
\usepackage{hyperref}
\ifdefined\darkmode
	\hypersetup{colorlinks=true,
		linkcolor=MaterialLightBlue300,
		citecolor=MaterialTeal300,
		urlcolor=MaterialLightBlue200}
\else
	\hypersetup{colorlinks=true,
		linkcolor=eoComment,
		citecolor=eoComment,
		urlcolor=eoComment}
\fi
\usepackage{ebproof}
\usepackage{array}

\usepackage{jbw-boxfigure}
\usepackage{jbw-fix-hyperref-autoref}

\usepackage{listings}
\usepackage[most]{tcolorbox}
\tcbuselibrary{listings}

%% ── Syntax colours ───────────────────────────────────────────────
\ifdefined\darkmode
	% Dark mode
	\definecolor{eoKeyword}{HTML}{f06292}%       commands   (pink)
	\definecolor{eoAttr}{HTML}{64b5f6}%          attributes (light blue)
	\definecolor{eoBuiltin}{HTML}{4dd0e1}%       builtins   (cyan)
	\definecolor{eoConst}{HTML}{b39ddb}%         types      (light purple)
	\definecolor{eoComment}{HTML}{9e9e9e}%       comments   (grey)
	\definecolor{eoString}{HTML}{ffb74d}%        strings    (light amber)
	\definecolor{eoPunct}{HTML}{bdbdbd}%         punctuation (light grey)
	\definecolor{eoFgMath}{HTML}{c8c0c8}%        foreground (softer)
\else
	% Light mode
	\definecolor{eoKeyword}{HTML}{e14775}%       commands   (pink)
	\definecolor{eoAttr}{HTML}{4a7fb5}%          attributes (muted blue)
	\definecolor{eoBuiltin}{HTML}{1c8ca8}%       builtins   (teal-cyan)
	\definecolor{eoConst}{HTML}{7058be}%         types      (purple)
	\definecolor{eoComment}{HTML}{3b5998}%       comments   (dark blue)
	\definecolor{eoString}{HTML}{cc7a0a}%        strings    (amber)
	\definecolor{eoPunct}{HTML}{918c8e}%         punctuation (grey)
	\definecolor{eoFgMath}{HTML}{3a3338}%        foreground (softer)
\fi

% Eunoia language definition for listings
\lstdefinelanguage{eunoia}{
	comment=[l]{;},
	string=[b]",
	morecomment=[s]{|}{|},% quoted symbols |...|
	alsoletter={-:@.},
	keywords={declare-const, declare-parameterized-const, declare-consts,
			declare-rule, declare-fun, declare-sort, declare-datatype,
			define, define-fun, define-const, define-sort,
			program, include, assume, assume-push, step, step-pop,
			let, match, as, par, forall, exists, lambda},
	keywords=[2]{:conclusion, :premises, :premise-list, :args,
			:requires, :assumption, :rule, :signature, :type,
			:right-assoc, :right-assoc-nil, :left-assoc, :left-assoc-nil,
			:chainable, :pairwise, :binder, :arg-list,
			:implicit, :opaque, :list, :sorry},
	% eo:: builtins handled via literate (see style below)
	keywords=[4]{Type, Proof},
	sensitive=true,
}

\lstdefinestyle{eo}{
language=eunoia,
basicstyle=\small\ttfamily\color{eoFgMath},
keywordstyle=\color{eoKeyword}\bfseries,
keywordstyle=[2]\color{eoAttr},
keywordstyle=[4]\color{eoConst},
commentstyle=\color{eoComment}\itshape,
stringstyle=\color{eoString},
showstringspaces=false,
columns=flexible,
keepspaces=true,
frame=none,
numbers=left,
numberstyle=\tiny\ttfamily\color{lineNoColor},
numbersep=6pt,
literate={(}{{\textcolor{eoPunct}{(}}}1
{)}{{\textcolor{eoPunct}{)}}}1
{[}{{\textcolor{eoPunct}{[}}}1
{]}{{\textcolor{eoPunct}{]}}}1
{eo::}{{\textcolor{eoBuiltin}{eo}\textcolor{eoPunct}{::}}}4
{$}{{\textcolor{eoConst}{\$}}}1,
	aboveskip=0.5\medskipamount,
	belowskip=0.5\medskipamount,
	}

	%% ── tcolorbox backgrounds ────────────────────────────────────────
	\ifdefined\darkmode
		\definecolor{eoBoxBg}{HTML}{2e3040}
		\definecolor{eoBoxFr}{HTML}{424454}
		\definecolor{lpBoxBg}{HTML}{2e3040}
		\definecolor{lpBoxFr}{HTML}{424454}
	\else
		\definecolor{eoBoxBg}{HTML}{fbfaf9}
		\definecolor{eoBoxFr}{HTML}{c5c1bf}
		\definecolor{lpBoxBg}{HTML}{fbfaf9}
		\definecolor{lpBoxFr}{HTML}{c5c1bf}
		\definecolor{lineNoColor}{HTML}{bfb9ba}
	\fi

	% tcolorbox style for Eunoia code blocks
	\newtcblisting{eolisting}{
		enhanced,
		listing only,
		listing options={style=eo, aboveskip=0pt, belowskip=0pt},
		colback=eoBoxBg,
		colframe=eoBoxFr,
		arc=0pt,
		boxrule=0.4pt,
		borderline west={1pt}{0pt}{eoBuiltin},
		left=14pt, right=2pt, top=2pt, bottom=2pt,
		before skip=\medskipamount,
		after skip=\medskipamount,
	}
	\newtcbinputlisting{\eofile}[1]{
		enhanced,
		listing only,
		listing file={#1},
		listing options={style=eo, aboveskip=0pt, belowskip=0pt},
		colback=eoBoxBg,
		colframe=eoBoxFr,
		arc=0pt,
		boxrule=0.4pt,
		borderline west={1pt}{0pt}{eoBuiltin},
		left=14pt, right=2pt, top=2pt, bottom=2pt,
		before skip=\medskipamount,
		after skip=\medskipamount,
		before upper=\listboxsetup,
	}

	% LambdaPi language definition and style
	\lstdefinelanguage{lambdapi}{
		comment=[l]{//},
		alsoletter={_},
		alsoother={@:$},
		keywords={symbol, rule, with, assert,
				require, open, type},
		keywords=[2]{constant, sequential, opaque, injective},
		keywords=[3]{TYPE},
		sensitive=true,
		mathescape=false,
		texcl=false,
	}

	\lstdefinestyle{lp}{
	language=lambdapi,
	basicstyle=\small\ttfamily\color{eoFgMath},
	keywordstyle=\color{eoKeyword}\bfseries,
	keywordstyle=[2]\color{eoAttr},
	keywordstyle=[3]\color{eoConst},
	commentstyle=\color{eoComment}\itshape,
	showstringspaces=false,
	columns=flexible,
	keepspaces=true,
	frame=none,
	numbers=left,
	numberstyle=\tiny\ttfamily\color{lineNoColor},
	numbersep=6pt,
	literate=
		{eo.Proof}{{\textcolor{eoBuiltin}{eo}\textcolor{eoPunct}{.}\textcolor{eoConst}{Proof}}}8
	{eo.app}{{\textcolor{eoBuiltin}{eo}\textcolor{eoPunct}{.}app}}6
	{eo.nil}{{\textcolor{eoBuiltin}{eo}\textcolor{eoPunct}{.}nil}}6
	{eo.}{{\textcolor{eoBuiltin}{eo}\textcolor{eoPunct}{.}}}3
	{eo::}{{\textcolor{eoBuiltin}{eo}\textcolor{eoPunct}{::}}}4
	{(}{{\textcolor{eoPunct}{(}}}1
	{)}{{\textcolor{eoPunct}{)}}}1
	{[}{{\textcolor{eoPunct}{[}}}1
	{]}{{\textcolor{eoPunct}{]}}}1
	{;}{{\textcolor{eoPunct}{;}}}1
	{:}{{\textcolor{eoPunct}{:}}}1
	{\{|}{{\textcolor{eoPunct}{\{|}}}2
	{|\}}{{\textcolor{eoPunct}{|\}}}}2
	{$}{{\textcolor{eoConst}{\$}}}1
{φ}{{$\varphi$}}1,
aboveskip=0.5\medskipamount,
belowskip=0.5\medskipamount,
}

% tcolorbox style for LambdaPi code blocks
\newtcblisting{lplisting}{
	enhanced,
	listing only,
	listing options={style=lp, aboveskip=0pt, belowskip=0pt},
	colback=lpBoxBg,
	colframe=lpBoxFr,
	arc=0pt,
	boxrule=0.4pt,
	borderline west={1pt}{0pt}{eoConst},
	left=14pt, right=2pt, top=2pt, bottom=2pt,
	before skip=\medskipamount,
	after skip=\medskipamount,
}
\newtcbinputlisting{\lpfile}[1]{
	enhanced,
	listing only,
	listing file={#1},
	listing options={style=lp, aboveskip=0pt, belowskip=0pt},
	colback=lpBoxBg,
	colframe=lpBoxFr,
	arc=0pt,
	boxrule=0.4pt,
	borderline west={1pt}{0pt}{eoConst},
	left=14pt, right=2pt, top=2pt, bottom=2pt,
	before skip=\medskipamount,
	after skip=\medskipamount,
	before upper=\listboxsetup,
}
% Thin separator line for use inside multi-part listing boxes
\newcommand{\listsep}{%
	\vspace{-\baselineskip}\vspace{7pt}%
	\noindent\kern-14pt\textcolor{eoBoxFr}{\rule{\dimexpr\linewidth+14pt+2pt}{0.3pt}}\par%
	\vspace{2pt}%
}
% Helper: zero out trivlist glue that listings inserts
\newcommand{\listboxsetup}{\topsep=0pt \partopsep=0pt \parsep=0pt \parskip=0pt}
% Multi-part listing boxes: use \listsep between \lstinputlisting calls
\newenvironment{eobox}{%
	\begin{tcolorbox}[
			enhanced, listing only,
			colback=eoBoxBg, colframe=eoBoxFr,
			arc=0pt, boxrule=0.4pt,
			borderline west={1pt}{0pt}{eoBuiltin},
			left=14pt, right=2pt, top=2pt, bottom=2pt,
			before skip=\medskipamount, after skip=\medskipamount,
		]\listboxsetup
		}{\end{tcolorbox}}
\newenvironment{lpbox}{%
	\begin{tcolorbox}[
			enhanced, listing only,
			colback=lpBoxBg, colframe=lpBoxFr,
			arc=0pt, boxrule=0.4pt,
			borderline west={1pt}{0pt}{eoConst},
			left=14pt, right=2pt, top=2pt, bottom=2pt,
			before skip=\medskipamount, after skip=\medskipamount,
		]\listboxsetup
		}{\end{tcolorbox}}
\lstset{defaultdialect=[]{eunoia}}

% Redefine example and definition to use italic headings
\let\example\relax
\let\endexample\relax
\spnewtheorem{example}{Example}{\itshape}{\rmfamily}
\def\exampleautorefname{example}
\let\definition\relax
\let\enddefinition\relax
\spnewtheorem{definition}{Definition}{\itshape}{\rmfamily}

\usepackage[backend=biber]{biblatex}
\addbibresource{inria.bib}

% small cases environment (brace + contents scaled to \small)
\newenvironment{scases}{%
	\left\{\,\small\begin{array}{@{}l@{\quad}l@{}}%
		}{%
	\end{array}\right.%
}

% font faces
\newcommand{\msf}{\mathsf}
\newcommand{\mbf}{\mathbf}
\newcommand{\mscr}{\mathscr}
\newcommand{\mtt}[1]{{\color{eoConst}{\texttt{#1}}}}% types/constants: Type, Bool, true, false, Proof
\newcommand{\kw}[1]{{\color{eoKeyword}{\texttt{\textbf{#1}}}}}% commands: declare-const, define, program, ...
\newcommand{\ev}[1]{{\color{eoFgMath}{\texttt{#1}}}}
\newcommand{\equo}{{\color{eoString}{\texttt{"}}}}% string quote
\newcommand{\estr}[1]{\equo{#1}\equo}% string literals

\newcommand{\cpc}{\msf{CPC}}
\newcommand{\tsm}[1]{{\small\text{#1}}}

\newcommand{\ptn}[1]{{\color{MaterialPurple500}{\texttt{#1}}}}
\newcommand{\mcomment}[1]{{\color{MaterialIndigo500}{\text{#1}}}}
\newcommand{\wrapp}[4]{{\color{#1}{#2}}{#4}{\color{#1}{#3}}}
\newcommand{\prn}[1]{\wrapp{MaterialGrey600}{\lparen}{\rparen}{#1}}

\newcommand{\cvc}[1]{\textsc{cvc{#1}}}
\newcommand{\verit}{\textsc{veriT}}
\newcommand{\eunoia}{\textsc{Eunoia}}

% ----- abstract syntax ----------
% plurality, binding
\newcommand{\plur}[2]{{#1}_{1}\ldots{#1}_{#2}}
\newcommand{\plurcomma}[2]{{#1}_{1}, \ldots, {#1}_{#2}}
\newcommand{\maybe}[1]{\wrapp{MaterialGrey600}{\langle}{\rangle_?}{#1}}
\newcommand{\many}[1]{\wrapp{MaterialGrey600}{\langle}{\rangle_+}{#1}}
\newcommand{\any}[1]{\wrapp{MaterialGrey600}{\langle}{\rangle_∗}{#1}}
\newcommand{\some}[1]{\wrapp{MaterialGrey600}{\langle}{\rangle}{#1}}
\newcommand{\binddot}[3]{{#1\,#2.\,#3}}

% binders
\newcommand{\lam}{\binddot{λ}}
\newcommand{\pii}{\binddot{Π}}
\newcommand{\alam}[1]{\binddot{\msf{λ}^{#1}}}
\newcommand{\apii}[1]{\binddot{\msf{Π}^{#1}}}
\newcommand{\lett}[2]{\msf{let}\ \prn{#1}\,\msf{in}\ #2}

% applications
\newcommand{\appSym}{\cdot}
\newcommand{\app}[2]{{#1}\appSym{#2}}
\newcommand{\apptwo}[3]{#1\,#2\,#3}
\newcommand{\appthree}[4]{#1\,#2\,#3\,#4}
\newcommand{\appfour}[5]{#1\,#2\,#3\,#4\,#5}
\newcommand{\appldots}[3]{\prn{\prn{\app{#1}{#2}}\,\ldots\,⋅\,{#3}}}

% universes
\newcommand{\PROP}{\mtt{PROP}}
\newcommand{\TYPE}{\mtt{TYPE}}
\newcommand{\KIND}{\mtt{KIND}}



% -------------- Eunoia ------------------------------------
% eunoia symbols
\newcommand{\epar}[1]{\wrapp{eoPunct}{\texttt{(}}{\texttt{)}}{#1}}
\newcommand{\eapp}[2]{\epar{#1\ #2}}
\newcommand{\earr}[1]{\epar{\ev{->}\ #1}}
\newcommand{\eprf}[1]{\epar{\mtt{Proof}\ #1}}
\newcommand{\ecase}[2]{\epar{#1\ #2}}
\newcommand{\eo}[1]{{\color{eoBuiltin}{\texttt{eo}}}{\color{eoPunct}{∷}}{\ev{#1}}}% builtins
\newcommand{\eotype}{\mtt{Type}}
\newcommand{\eoreqs}[3]{\eapp{\eo{requires}}{#1\ #2\ #3}}
\newcommand{\eoquote}[1]{\eapp{\eo{quote}}{#1}}

\newcommand{\attr}[1]{{\color{eoAttr}{\texttt{:#1}}}}% attributes
\newcommand{\premises}[1]{\attr{premises}\,\epar{#1}}
\newcommand{\premiselist}[2]{\attr{premise-list}\ #1\ #2}
\newcommand{\requires}[1]{\attr{requires}\ \epar{#1}}
\newcommand{\conclusion}[1]{
	\attr{conclusion}\ #1
}


\newcommand{\sym}[1]{\msf{s}_{\msf{#1}}}
% \newcommand{\eoapp}[2]{\prn{#1\ #2}}
% \newcommand{\earr}[1]{\mtt{#1}}
\newcommand{\ctyp}{\mbf{ctyp}}
\newcommand{\att}{\mbf{att}}

\newcommand{\prf}{\mbf{prf}}
\newcommand{\step}{\mbf{step}}

\newcommand{\eor}{\texttt{or}}
\newcommand{\econcat}[3]{\eapp{\eo{list\_concat}}{#1\ #2\ #3}}
\newcommand{\econs}[3]{\eapp{#1}{#2\ #3}}
\newcommand{\false}{\texttt{false}}

\newcommand{\rcn}[1]{\attr{right-assoc-nil}\ {#1}}
\newcommand{\rc}{\attr{right-assoc}}

\newcommand{\lcn}[1]{\attr{left-assoc-nil}\ {#1}}
\newcommand{\lc}{\attr{left-assoc}}

\newcommand{\ch}[1]{\attr{chainable}\ {#1}}
\newcommand{\pw}[1]{\attr{pairwise}\ {#1}}
\newcommand{\bd}[1]{\attr{binder}\ {#1}}
\newcommand{\al}[1]{\attr{arg-list}\ {#1}}
\newcommand{\eovar}[2]{\eapp{\eo{var}}{#1\ #2}}

% --- builtin operators (sec:eo-builtin) ---
\newcommand{\eoBuiltinFig}{
	\begin{array}[t]{r@{\ }l@{\ }l@{\quad}l}
		β & {⦂}
		  & \eo{ite}\ ∣\ \eo{eq}\ ∣\ \eo{requires}\ ∣\ \eo{typeof}\ ∣\ \eo{var}\ ∣\ \ldots
		  & \mcomment{(core)}
		\\
		  & {∣}
		  & \eo{and}\ ∣\ \eo{or}\ ∣\ \eo{not}\ ∣\ \eo{xor}
		  & \mcomment{(boolean)}
		\\
		  & {∣}
		  & \eo{add}\ ∣\ \eo{mul}\ ∣\ \eo{neg}\ ∣\ \eo{zdiv}\ ∣\ \ldots
		  & \mcomment{(arithmetic)}
		\\
		  & {∣}
		  & \eo{nil}\ ∣\ \eo{cons}\ ∣\ \eo{list\_concat}\ ∣\ \eo{list\_len}\ ∣\ \ldots
		  & \mcomment{(list)}
	\end{array}
}


\newcommand{\cons}{{\text{cons}}}
\newcommand{\nil}{{\text{nil}}}
\newcommand{\foldl}{\mbf{foldl}}
\newcommand{\foldr}{\mbf{foldr}}
\newcommand{\elab}[1]{\mbf{elab}_{\prn{#1}}}
\newcommand{\glue}[1]{\mbf{glue}_{\prn{#1}}}
\newcommand{\bool}{\texttt{Bool}}
% ---- eo commands ------
\newcommand{\dc}[2]{\epar{\kw{declare-const}\ {#1}\ {#2}}}
\newcommand{\dca}[3]{\epar{\kw{declare-const}\ {#1}\ {#2}\ {#3}}}

% --- figures for `eoas.tex` ------------------------------
% s & {⦂}\ \sym 0 ∣ \sym 1 ∣ \ldots
% 		& \mcomment{(symbols)}         &
%
\newcommand{\eoTermSyntax}{
	\begin{array}[t]{r@{\ }l@{\quad}l@{\qquad}r@{\ }l@{\quad}l}
		%!----------------------------------------!%
		t & {⦂}\ s \mid ℓ \mid \eapp{s}{\vec t} \mid \eapp{s}{\epar{\vec{ρ}}\ t} & \mcomment{(terms)}      &
		ρ & {⦂}\ \epar{s\ t\ \maybe{ν}}                                          & \mcomment{(parameters)}   \\[1mm]
		ℓ & {⦂}\ \msf{int}(n) ∣ \msf{str}(s)                                     & \mcomment{(literals)}   &
		c & {⦂}\ \epar{t\ t'}                                                    & \mcomment{(cases)}        \\[1mm]
	\end{array}
}

\newcommand{\eoAttrSyntax}{
	\begin{array}[t]{r@{\ }c@{\ }l@{\quad}l}
		ν & {⦂} & \attr{implicit} \mid \attr{list}               & \mcomment{(prm. attributes)}   \\
		α & {⦂} & \attr{right-assoc} ∣ \attr{right-assoc-nil}\ t & \mcomment{(const. attributes)} \\
		  & {∣} & \attr{left-assoc} ∣ \attr{left-assoc-nil}\ t                                    \\
		  & {∣} & \attr{chainable}\ s ∣ \attr{pairwise}\ s                                        \\
		  & {∣} & \attr{binder}\ s ∣ \attr{arg-list}\ s
	\end{array}
}

\newcommand{\eoCommandSyntax}{
	\begin{array}[t]{r@{\ }l@{\quad}l}
		δ & \begin{array}[t]{c@{\ }l}
			    {⦂} & \epar{\kw{declare-const}\ s\ t\ \maybe{α}}                               \\
			    {∣} & \epar{\kw{declare-parameterized-const}\ s\ \epar{\vec{ρ}}\ t\ \maybe{α}} \\
			    {∣} & \epar{\kw{declare-consts}\ ℓ\ t}                                         \\
			    {∣} & \epar{\kw{define}\ {s} \ \epar{\vec ρ} \ t\ \ \maybe{\attr{type}\,t'}}   \\
			    {∣} & \epar{\kw{program}\ {s} \ \epar{\vec ρ}\
			    \attr{signature}\,\epar{\vec{t}}\,t'\
			    \epar{\vec{c}}
			    }                                                                              \\
			    {∣} & {\color{eoPunct}{\texttt{(}}}
			    \kw{declare-rule}\ s\ \epar{\vec ρ}\
			    \\ & \quad \maybe{\attr{assumption}\ t_{\text{assm}}}\
			    \\ & \quad \maybe{
			    \attr{premises}\ \epar{\vec{t}_{\text{prem}}}
			    \mid
			    \attr{premise-list}\ t_{\text{prem}}\ t_{\text{op}}
			    }\
			    \\ & \quad \maybe{\attr{args}\ \epar{\vec{t}_{\text{args}}}}\
			    \\ & \quad \maybe{\attr{requires}\ \epar{\vec{c}}}\
			    \\ & \quad \attr{conclusion}\ {t_{\text{conc}}}
			    {\color{eoPunct}{\texttt{)}}}                                                  \\
			    {∣} & \epar{\kw{include}\ μ}
		    \end{array}
		  & \mcomment{(std. commands)}
	\end{array}
}

\newcommand{\eoPrfCommandSyntax}{
	\begin{array}[t]{r@{\ }l@{\quad}l}
		π & \begin{array}[t]{c@{\ }l}
			    {⦂} & \epar{\kw{assume}\ s\ φ}      \\
			    {∣} & \epar{\kw{assume-push}\ s\ φ} \\
			    {∣} & \epar{\kw{step}\ s\ φ\
			    \attr{rule}\ s'\
			    \maybe{\attr{premises}\ {\vec ψ}}\
			    \maybe{\attr{args}\ {\vec t}}}      \\
			    {∣} & \epar{\kw{step-pop}\ s\ φ\
			    \attr{rule}\ s'\
			    \maybe{\attr{premises}\ {\vec ψ}}\
			    \maybe{\attr{args}\ {\vec t}}}      \\
		    \end{array}
		  & \mcomment{(prf. commands)}
	\end{array}
}

%
%
% --- macros for `lp.tex` ------------------------------
\newcommand{\esc}[1]{\wrapp{MaterialGrey600}{⦃}{⦄}{#1}}
\newcommand{\var}[1]{\msf{v}_{#1}}
\newcommand{\lbl}[2]{\rname{#1}\quad{#2}}
\newcommand{\lpTermSyntax}{
	\begin{array}[t]{r@{\ }l@{\quad}r}
		% ------------------------------------- %
		μ & {⦂}\ \TYPE ∣ \KIND
		  & \mcomment{(universes)}                                                             \\[1mm]
		% ------------------------------------- %
		t & {⦂}\ x ∣ κ ∣ μ ∣ \prn{\app{t}{t'}} ∣ \prn{\lam{x : t}{t'}} ∣ \prn{\pii{x : t}{t'}}
		  & \mcomment{(terms)}
		% ------------------------------------- %
	\end{array}
}

\newcommand{\lpCtxSyntax}{
	\begin{array}[t]{r@{\ }l@{\quad}r}
		% ------------------------------------- %
		γ & {⦂}\ \prn{x : t} \mid \prn{t ↪ t'} \
		  & \mcomment{(context elements)}        \\
		% ------------------------------------- %
		Γ & {⦂}\ \any{γ}
		  & \mcomment{(contexts)}
		% ------------------------------------- %
	\end{array}
}

\newcommand{\judge}[3]{{#1} ⊢_{Σ} {#2} : {#3}}
\newcommand{\wf}[1]{\msf{wf}\,{#1}}
% typing rules
\newcommand{\typeRule}{
	\begin{prooftree}
		\hypo{\wf Γ}
		\infer1{ \judge Γ \TYPE \KIND }
	\end{prooftree}
}

\newcommand{\varRule}{
	\begin{prooftree}
		\hypo{ \wf Γ}
		\infer1[$\prn{x : t}∈ Γ$]{ \judge Γ x t }
	\end{prooftree}
}

\newcommand{\constRule}{
	\begin{prooftree}
		\hypo{\wf Γ}
		\hypo{\judge {} t μ}
		\infer2[$\prn{κ : t}∈ Σ$]{ \judge Γ κ t }
	\end{prooftree}
}

\newcommand{\prodRule}{
	\begin{prooftree}
		\hypo{\judge Γ t {\TYPE}}
		\hypo{\judge{Γ, \prn{x : t}}{t'}{μ'}}
		\infer2{\judge Γ {\prn{\pii {x : t} {t'}}} {μ'} }
	\end{prooftree}
}

\newcommand{\absRule}{
	\begin{prooftree}
		\hypo{\judge{Γ, \prn{x : t}} e {t'}}
		\hypo{\judge Γ {\prn{\pii {x : t} {t'}}} μ}
		\infer2{\judge Γ {\prn{\lam {x : t} {e}}} {\prn{\pii {x : t} {t'}}} }
	\end{prooftree}
}

\newcommand{\appRule}{
	\begin{prooftree}
		\hypo{\judge Γ e {\prn{\pii{x : t}{t'}}}}
		\hypo{\judge Γ {e'} t}
		\infer2{\judge Γ {\prn{\app{e}{e'}}} {t'[x ↦ {e'}]}}
	\end{prooftree}
}

\newcommand{\convRule}{
	\begin{prooftree}
		\hypo{\judge Γ e t}
		\hypo{\judge Γ {t'} {μ}}
		\infer2[$\prn{t ≡_{βΣ} t'}$]{\judge Γ {e} {t'}}
	\end{prooftree}
}

\newcommand{\wfRuleEmp}{
	\begin{prooftree}
		\infer0{ \wf{\msf{∅}} }
	\end{prooftree}
}

\newcommand{\wfRuleVar}{
	\begin{prooftree}
		\hypo{ \judge Γ t μ }
		\infer1[$x ∉ \mbf{dom}(Γ)$]{ \wf{(Γ, \prn{x : t})}}
	\end{prooftree}
}

\newcommand{\rname}[1]{\mcomment{($\textsc{#1}$)}}

\newcommand{\lpTypingRules}{
	\begin{array}[t]{c}
		\lbl{wf0}{\wfRuleEmp} \qquad \lbl{wf+}{\wfRuleVar}
		\\[6mm]
		\lbl{var}{\varRule}    \qquad  \lbl{con}{\constRule}
		\\[6mm]
		\lbl{univ}{\typeRule}  \qquad \lbl{prod}{\prodRule}
		\\[6mm]
		\lbl{fun}{\absRule} \\[6mm]
		\lbl{app}{\appRule} \\[6mm]
		\lbl{conv}{\convRule}
	\end{array}
}



% --- concrete syntax for lambdapi
\newcommand{\lambdapi}{\textsc{LambdaPi}}
\newcommand{\llam}{\binddot{\msf{λ}}}
\newcommand{\ppii}{\binddot{\msf{Π}}}

% \newcommand{\symbol}{\mtt{symbol}}

\newcommand{\prm}[1]{\mbf{#1}}
\newcommand{\xsym}[1]{\mtt{@}{#1}}
\newcommand{\meta}[1]{\ptn{?{#1}}}

\newcommand{\ptrn}[1]{\ptn{\$#1}}
\newcommand{\pv}{\mbf{pvars}}

\newcommand{\imp}[1]{[#1]}
\newcommand{\xpl}[1]{\left(#1\right)}

\newcommand{\lpTermsConcrete}{
	\begin{array}[t]{r@{\ }l@{\quad}r}
		t & {⦂}\ s ∣ \imp{t} ∣ \prn{\app{t}{t'}}
		∣ \prn{\llam{ρ}{t}} ∣ \prn{\ppii{ρ}{t}}
		  & \mcomment{(terms)}
	\end{array}
}

\newcommand{\lpPatternSyntax}{
	\begin{array}[t]{r@{\ }c@{\ }l@{\quad}r@{\qquad}r@{\ }c@{\ }l@{\quad}r}
		m & {⦂}
		  & \mtt{constant} ∣ \mtt{sequential}
		  & \mcomment{(modifiers)}
		  &
		θ & {⦂}
		  & \ptrn{x} ∣ s\,\any{θ}
		  & \mcomment{(patterns)}
		\\[1mm]
		ρ & {⦂}
		  & \xpl{s : t} ∣ \imp{s : t}
		  & \mcomment{(parameters)}
		  &
		% ------------------------------------- %
		r & {⦂}
		  & \prn{{s\,\any{θ}} ↪ {θ'}}
		  & \mcomment{(rewrite rules)}
		% ------------------------------------- %
	\end{array}
}

\newcommand{\lpConcreteSyntax}{
	\begin{array}[t]{r@{\ }c@{\ }l@{\quad}r}
		c & {⦂}
		  & \maybe{m}\, \mtt{symbol}\ s \ \any{ρ} \,{: t} \ \maybe{≔ t'};
		  & \mcomment{(commands)}
		\\
		  & {∣}
		  & \mtt{rule}\ r\ \any{\mtt{with}\ r'};
		\\
		  & {∣}
		  & \mtt{require}\,\mtt{open}\,\many{μ};
	\end{array}}

% \newcommand{\symb}[1]{\mtt{symbol}\ {#1};}

% \newcommand{\symb}[1]{\mtt{rule}\ {#1};}

\newcommand{\Set}{\msf{Set}}
\newcommand{\El}[1]{\msf{τ}\,{#1}}

% ------- translation stuff -----------
\newcommand{\Term}{\msf{Term}}

%
\newcommand{\EO}{\msf{EO}}
\newcommand{\LP}{\msf{LP}}
\newcommand{\trans}[1]{\left⟦{#1}\right⟧}
\newcommand{\tm}[1]{\trans{#1}_\msf{tm}}
\newcommand{\ty}[1]{\trans{#1}_\msf{ty}}
\newcommand{\cmd}[1]{\trans{#1}_\msf{cmd}}
\newcommand{\ctx}[1]{\trans{#1}_\msf{ctx}}


\begin{document}
\section{Eunoia}\label{sec:eunoia}
% fig:eo-syntax is in intro.tex so it floats to the top of page 2
% ---------------------------------------------------------
This section presents the fragment of Eunoia targeted by our
translation, together with the elaboration step that resolves
$n$-ary applications to normal form.
%
Eunoia currently has no published formal semantics:
its meaning is defined operationally by Ethos.
%
The presentation here was informed by the
Ethos user manual~\cite{EthosUser_manualmdMain},
the Ethos source, and through discussions with the
\cvc{5}/Eunoia developers.
%
% We accordingly do not give a typing relation or
% algorithm for Eunoia, nor a model-theoretic semantics
% analogous to the one in SMT-LIB~2.6~\cite{Barrett2015-standard}.


% ---------------------------------------------------------
\subsection{Syntax, Commands, and Signatures}\label{sec:eo-sig}
% ---------------------------------------------------------
%
\Autoref{fig:eo-syntax} defines the concrete syntax for Eunoia,
where the metavariable $s$ ranges over a set of \emph{symbols}.
%
A term is either a symbol $s$, a literal $ℓ$
(integer, rational, or string), an \emph{application}
$\eapp{s}{\vec t}$, or a \emph{binder application}
$\eapp{s}{\epar{\vec ρ}\ t}$.
%
A \emph{parameter} $\epar{s\ t\ \maybe{ν}}$ declares a symbol
$s$ of type $t$ with an optional \emph{parameter attribute} $ν$
% ($\attr{implicit}$ marks $ρ$ as
% inferable from later arguments, and $\attr{list}$
% marks it as an $f$-list pattern (see below).
Similarly, a symbol may have a \emph{constant attribute} $α$
that specifies how its $n$-ary applications are elaborated
to binary form (see \autoref{sec:eo-elab}).
%
Commands are divided into \emph{standard commands}
and \emph{proof commands}.

A \emph{signature} $Σ$ is a set of \emph{declarations},
where each declaration is either a \emph{typing} $s\prn{\vec ρ} : t$,
an \emph{attribution} $s\prn{\vec ρ} ∷ α$,
or a \emph{rewrite rule} $ℓ ↪ r$ with parameters $\vec \rho$.
%
\autoref{tab:cmd-decls} maps each command to
the set of declarations it contributes.
%
Any sequence of standard commands $\vec δ$ thus
determines a signature: the union of the
declarations contributed by each command.
%
A \emph{proof script} $Υ = \prn{\vec δ ⨾ \vec π}$
pairs a sequence of standard commands $\vec δ$
(the \emph{preamble}) with a sequence of proof
commands $\vec π$ (the \emph{body}); $Υ$ likewise
determines a signature, which we call a \emph{proof}.
%

% ---------------------------------------------------------
\begin{boxfigure}[t!]{fig:eo-builtins}
	{Eunoia built-in symbols, grouped by family.}
	\alttext{Built-in symbols grouped by family---core (eo::eq, eo::requires,
		eo::quote, eo::var), boolean (eo::and, eo::or, eo::not, eo::xor),
		arithmetic (eo::add, eo::mul, eo::neg, eo::zdiv), and
		list (eo::nil, eo::cons, eo::list\_concat, eo::list\_len).}{%
		\vspace{4pt}\centerline{$\eoBuiltinFig$}\vspace{4pt}}
\end{boxfigure}
% ---------------------------------------------------------
%
Eunoia also provides a set of \emph{built-in symbols}~(\autoref{fig:eo-builtins})
that operate on the syntactic structure of terms.
%
The symbols $\eo{nil}$, $\eo{cons}$, $\eo{list\_concat}$
construct and combine \emph{$f$-lists}
(i.e., nested applications of a binary $f$,
defined formally in \autoref{sec:eo-elab});
$\eoreqs{t_1}{t_2}{t_3}$ acts as a guard, evaluating
to $t_3$ only when $t_1$ and $t_2$ are syntactically
identical; $\eoquote{t}$ marks $t$ as a \emph{quoted
	argument} whose value may appear in the return type
of a rule or program (Eunoia's mechanism for
object-level dependent types); and $\eo{var}$
constructs variable-like terms used by quantifiers
(\autoref{sec:eo-elab}, item~2).


\begin{table}[t!]
	\centering
	\caption{Declarations produced by each Eunoia command.
		Parts in $\maybe{\cdot}$ are present only when the
		corresponding attribute is supplied. Each program case
		$c_i = \ecase{l_i}{r_i}$ contributes one rewrite rule.}
	\label{tab:cmd-decls}
	\small
	\begin{tabular}{@{}lp{0.46\linewidth}@{}}
		\hline
		\textbf{Command} & \textbf{Declarations}                            \\
		\hline
		$\epar{\kw{declare-const}\ s\ t\ \maybe{α}}$
		                 & $s : t,\,\maybe{s ∷ α}$                          \\
		$\epar{\kw{declare-parameterized-const}\ s\ \epar{\vec ρ}\ t\ \maybe{α}}$
		                 & $s\prn{\vec ρ} : t,\,\maybe{s ∷ α}$              \\
		$\epar{\kw{declare-consts}\ ℓ\ t}$
		                 & binds the literal class $ℓ$ to type $t$          \\
		$\epar{\kw{define}\ s\ \epar{\vec ρ}\ t\ \maybe{\attr{type}\ t'}}$
		                 & $s\prn{\vec ρ} ↪ t,\,\maybe{s\prn{\vec ρ} : t'}$ \\
		$\epar{\kw{program}\ s\ \epar{\vec ρ}\ \attr{signature}\ \epar{\vec t}\ t'\ \epar{c_1 \ldots c_m}}$
		                 & $s\prn{\vec ρ} : \earr{\vec t\ t'}$\,;\
		$l_i ↪ r_i$ for each $c_i = \ecase{l_i}{r_i}$                       \\
		$\epar{\kw{declare-rule}\ s\ \epar{\vec ρ}\ \ldots}$
		                 & $s\prn{\vec ρ} : τ^★$
		(\autoref{sec:eo-rules})                                            \\
		$\epar{\kw{include}\ μ}$
		                 & declarations of module~$μ$                       \\
		$\epar{\kw{assume}\ s\ φ}$
		\,/\,$\epar{\kw{assume-push}\ s\ φ}$
		                 & $s : \eprf{φ}$                                   \\
		$\epar{\kw{step}\ s\ φ\ \attr{rule}\ r\ \maybe{\attr{args}\ \epar{\vec t}}\
				\maybe{\attr{premises}\ \epar{\vec p}}}$
		                 & $s : \eprf{φ}$,\,
		$s ↪ \eapp{r}{\vec t\ \vec p}$                                      \\
		\hline
	\end{tabular}
\end{table}

\subsection{Rule Declarations and Proofs.}\label{sec:eo-rules}
%
Given an inference rule declaration
$\epar{\kw{declare-rule}\ s\ \epar{\vec ρ}\ \ldots\ \conclusion φ}$,
the derivation of its contributed type declaration
$s\prn{\vec\rho} : \tau^\star$ is given thus:
%
\begin{enumerate}
	\item Given
	      $\requires{\ecase{l_1}{r_1}\ldots\ecase{l_k}{r_k}}$,
	      define $φ^★$ below. Otherwise, $φ^★ ≔ φ$.
	      $$φ^★ ≔ \eoreqs{l_1}{r_1}{
			      \epar{
				      \,\ldots\,\eoreqs{l_k}{r_k}{φ}}
		      }
	      $$

	\item If the declaration provides no premises,
	      let $τ^★ ≔ \eprf{φ^★}$.
	      Otherwise, define $τ^★$ thus:
	      \begin{enumerate}
		      \item Given
		            $\attr{premises}\ \epar{\plur ψ m}$,
		            let:
		            $$τ^★ ≔ \earr{\eprf{ψ_1}\,\ldots\,\eprf{ψ_m}\ \eprf{φ^★}}$$
		            %
		      \item Given $\premiselist ψ f$ for
		            some associative symbol $f$, redefine:
		            $$φ^★ ≔\eoreqs{\eapp{\eo{is\_list}}{f\ ψ}}{\ev{true}}{φ^★}$$
		            and let $τ^★ ≔ \earr{\eprf{ψ}\ \eprf{φ^★}}$.
	      \end{enumerate}

	\item If the declaration provides
	      $\attr{args}\ \epar{\plur t n}$,
	      then the type contributed for $s$ is:
	      $$s\prn{\vec ρ}
		      : \earr{\eoquote{t_1}\,\ldots\,\eoquote{t_n}\ τ^★}
	      $$
	      Otherwise, $s\prn{\vec ρ}: τ^★$.
\end{enumerate}

\begin{example}\label{ex:eo-rules}
	Consider two rule declarations from
	\texttt{cpc/rules/Booleans.eo}.
	%
	\begin{eobox}
		\lstinputlisting[style=eo, aboveskip=0pt, belowskip=0pt]{listings/rule-modus-ponens.eo}%
		\lstinputlisting[style=eo, aboveskip=0pt, belowskip=0pt]{listings/rule-cnf-implies-pos.eo}%
	\end{eobox}
	%
	\noindent
	The $\ev{modus\_ponens}$ rule has only $\attr{premises}$ and
	$\attr{conclusion}$, so item~2a gives:
	$$
		\ev{modus\_ponens}(\vec ρ)  :
		{\color{eoPunct}{\texttt{(}}}\ev{->}\
		\eprf{\ev{F1}}\
		\eprf{\eapp{\ev{=>}}{\ev{F1}\ \ev{F2}}}\
		\eprf{\ev{F2}}
		{\color{eoPunct}{\texttt{)}}}
	$$
	%
	The $\ev{cnf\_implies\_pos}$ rule has
	$\attr{args}\ \epar{\eapp{\ev{=>}}{\ev{F1}\ \ev{F2}}}$;
	item~3 makes this an
	$\eo{quote}$-typed parameter that pattern-matches at
	the point of use, where $φ$ is the term given by
	its $\attr{conclusion}$ attribute:
	$$
		\ev{cnf\_implies\_pos}(\vec ρ)  :
		{\color{eoPunct}{\texttt{(}}}\ev{->}\
		\eoquote{\eapp{\ev{=>}}{\ev{F1}\ \ev{F2}}}\
		\eprf{φ}
		{\color{eoPunct}{\texttt{)}}}
	$$
	% TODO. re-add `reordering` and and_intro` to listing and prose.
	Other rules combine these attributes:
	$\ev{reordering}\ \attr{args}\,\attr{requires}$
	uses item~3 plus a guard via item~1, and
	$\ev{and\_intro}\ \premiselist{\ev F}{\ev{and}}$
	guards $φ$ with $\eapp{\eo{is\_list}}{\ev{and}\ \ev{F}}$
	via item~2b.
\end{example}


% ---------------------------------------------------------
\subsection{Elaboration}\label{sec:eo-elab}
% ---------------------------------------------------------
%
Unlike standard LF-style frameworks, Eunoia supports
$n$-ary applications of binary symbols, with the
unrolling determined by a \emph{constant attribute}
$α$ (\autoref{fig:eo-syntax}) assigned to the symbol
at its declaration.
%
After elaboration, every application of a declared
constant $s$ uses the built-in higher-order
application operator $\mtt{\_}$, so that
$\eapp{s}{t_1\ t_2}$ becomes $\eho{\eho{s}{t_1}}{t_2}$;
only built-in and program symbols are applied
directly.
%
For example, given $f$ with attribute $\rcn{t_\nil}$,
the \emph{$f$-list} with elements $\plur t n$ is the
nested application
$\eho{\eho{f}{t_1}}{\eho{\eho{f}{\ldots}}{\eho{\eho{f}{t_n}}{t_\nil}}}$.
%
Similarly, $\bd{t_\cons}$ specifies that a binder
application $\eapp{s}{\epar{\vec ρ}\ t}$ is elaborated
so that each parameter $\epar{x_i\ t_i}$ in $\vec ρ$
becomes a term $\epar{\eo{var}\ \estr{x_i}\ t_i}$
within a $t_\cons$-list.

The $\mbf{elab}$ operator below captures these
elaboration strategies uniformly.

\newcommand{\body}{\text{body}}
\newcommand{\vars}{\text{vars}}
\begin{definition}\label{def:elab}
	For any $Δ$ and $\vec ρ$, let $\elab{Δ,\vec ρ}$
	be the least operator such that:
	\begin{enumerate}
		\item For any symbol $s$, $\elab{Δ,\vec ρ}[s] = s$.
		      %
		\item For a binder application
		      $\eapp{s}{\epar{\vec v}\ t}$ with
		      $Δ ⊢ s ∷ \bd t_\cons$, each parameter
		      $\epar{x_i\ t_i} ∈ \vec v$ is reified as
		      $X_i ≔ \epar{\eo{var}\ \estr{x_i}\ t_i}$,
		      and:
		      $$ \elab{Δ,\vec ρ}[\eapp{s}{\epar{\vec v}\ t}]
			      =
			      \eapp{s}{t_\vars\ t_\body}
			      \quad\text{where}\quad
			      \begin{array}{l@{\ {≔}\ }l}
				      t_\vars &
				      \elab{Δ, \vec ρ}[\eapp{t_\cons}{\plur X n}]
				      \\
				      t_\body &
				      \elab{Δ, \vec ρ}
				      [\mbf{subst}_{\prn{\vec x, \vec X}}[t]]
			      \end{array}
		      $$

		      %
		\item For an application $\eapp s {\vec t}$,
		      let $t_i' ≔ \elab{Δ, \vec ρ}[t_i]$ in:
		      \begin{enumerate}
			      \item If $s$ is a built-in or
			            $s ∈ Δ.\msf{prog}$, then
			            $\elab{Δ, \vec ρ}[\eapp{s}{\vec t}] =
				            \eapp{s}{\vec t'}$.
			      \item If $s ∈ \vec ρ$, then
			            $\elab{Δ, \vec ρ}[\eapp{s}{\vec t}] =
				            \eho{\ldots\eho{\eho{s}{{t_1'}}}{\ldots}}{{t_n'}}$.

			      \item If $Δ ⊢ s ∷ α$ for an associative
			            $α$, the result is one of:
			            $$
				            \begin{scases}
					            \foldr(G, t_\nil, \vec t')
					            & \tsm{$\prn{α = \rcn{t_\nil}}$}
					            \\
					            \foldl(G, t_\nil, \vec t')
					            & \tsm{$\prn{α = \lcn{t_\nil}}$}
					            \\
					            \foldr(G, t'_n, t'_1\ldots t'_{n-1})
					            & \tsm{$\prn{α = \rc}$}
					            \\
					            \foldl(G, t'_1, t'_2\ldots t'_{n})
					            & \tsm{$\prn{α = \lc}$}
				            \end{scases}
			            $$
			            where the combining function
			            $G(x,y)$ is $\econcat{s}{x}{y}$ if
			            $\vec ρ ⊢ x ∷ \attr{list}$ and
			            $\eho{\eho{s}{x}}{y}$ otherwise.

			      \item If $Δ ⊢ s ∷ α$ for a non-associative
			            $α$, then
			            $\elab{Δ,\vec ρ}[\eapp{s}{\vec t}] =
				            \elab{Δ,\vec ρ}[\eapp{t_\cons}{\vec u}]$
			            where:
			            $$
				            \vec u ≔
				            \begin{scases}
					            \eapp{s}{t_1 t_2}\,\eapp{s}{t_2 t_3}
					            \ldots\,\eapp{s}{t_{n-1} t_n}
					            &
					            \tsm{$\prn{α = \ch {t_\cons}}$}
					            \\
					            \eapp{s}{t_i t_j} \text{ for all } i < j
					            &
					            \tsm{$\prn{α = \pw {t_\cons}}$}
					            \\
					            \vec t
					            &
					            \tsm{$\prn{α = \al {t_\cons}}$}
				            \end{scases}
			            $$

		      \end{enumerate}
	\end{enumerate}
\end{definition}

\begin{example}\label{ex:elab}
	Since $Δ ⊢ \ev{or} ∷ \rcn\false$, item~3c applies
	with $G(x,y) = \eho{\eho{\ev{or}}{x}}{y}$ and the
	right fold gives:
	$$ \elab{Δ, \vec ρ}[\eapp{\ev{or}}{\ev x\ \ev y\ \ev z}]
		=
		\eho{\eho{\ev{or}}{\ev x}}{
			\eho{\eho{\ev{or}}{\ev y}}{
				\eho{\eho{\ev{or}}{\ev z}}{\false}}}
	$$
	%
	For $Δ ⊢ \ev{=} ∷ \ch{\ev{and}}$, item~3d
	first chains the equalities and then recurses
	through $\ev{and}$'s $\rcn\false$ attribute:
	$$
		\elab{Δ, \vec ρ}[\eapp{\ev{=}}{\ev a\ \ev b\ \ev c\ \ev d}]
		=
		\elab{Δ, \vec ρ}[\eapp{\ev{and}}{
				\eapp{\ev{=}}{\ev a\ \ev b}\
				\eapp{\ev{=}}{\ev b\ \ev c}\
				\eapp{\ev{=}}{\ev c\ \ev d}}]
	$$
\end{example}

\end{document}
